\section{Practical Implications}
\label{sec:practical-implications}

The theoretical and experimental results of this paper have several practical
consequences for the use of stability diagnostics in machine learning.

\subsection{Cross-Validation Stability Is Not Adversarial Robustness}

Leave-one-out cross-validation is a standard tool for model assessment. When a
model has low LOO error variance, practitioners often infer that the model is
``stable.'' Our results show that this inference is unwarranted: a model can be
perfectly LOO-stable while simultaneously being highly vulnerable to a single
adversarially chosen replacement. The LOO/replace-one gap means that
cross-validation stability certificates do not automatically transfer to
adversarial replacement scenarios.

This is not a theoretical curiosity. The experimental results of
\Cref{sec:tabular-benchmark} show that the gap appears on standard UCI datasets
with measurable frequency. A practitioner evaluating a $k$-NN classifier on the
Digits dataset with $k=1$ would see zero LOO instability in every replicate, yet
a single well-chosen replacement flips the prediction in every case.

\subsection{Margin-Based Vulnerability Detection Is Practical}

The StabilityGapDiagnostic algorithm (\Cref{sec:diagnostic-algorithm}) provides
a computationally efficient way to detect the gap. The margin criterion
$|M_k(S, x)| \le 2$ requires only the computation of the signed vote margin for
each query, which is $O(n \log n)$ for a single query and straightforward to
compute for all queries in a dataset.

The diagnostic outputs:
\begin{itemize}
    \item A list of {\em vulnerable queries} where the margin is within the
        replace-one flip range.
    \item {\em Margin-crossing certificates} that demonstrate explicit
        replacement perturbations capable of changing the prediction.
    \item An overall {\em LOO/replace-one gap score} that quantifies the
        discrepancy between the two stability notions.
\end{itemize}

\subsection{Guidelines for Practitioners}

Based on the analysis, we offer the following practical recommendations:

\begin{enumerate}[leftmargin=*,label=(\arabic*)]
    \item \textbf{Do not equate LOO stability with replacement robustness.}
        A low LOO error is not evidence that the model is insensitive to
        single-point adversarial perturbations. Always check both LOO and
        replace-one indicators when robustness is a concern.
    \item \textbf{Monitor the margin distribution.} The signed margin
        $M_k(S, x)$ is a more informative diagnostic than aggregate loss.
        Queries with margin $|M_k| \le 2$ are potentially vulnerable to a
        single replacement.
    \item \textbf{Avoid small $k$ in high-stakes settings.} The LOO/replace-one
        gap is most frequent at small $k$ and in high-dimensional spaces where
        near-distance ties are common. If the environment is adversarial,
        increasing $k$ reduces the vulnerability.
    \item \textbf{Use the diagnostic before deployment.} Run
        \textsc{StabilityGapDiagnostic} on the training set to detect
        vulnerable queries. If the gap exceeds the application's tolerance,
        consider switching to a smoother classifier or using an ensemble.
\end{enumerate}

\subsection{Limitations of the Practical Analysis}

The diagnostic is designed for deterministic $k$-NN on finite metric spaces.
Extending it to randomized classifiers, real-valued scores, or non-metric
similarity measures requires additional theoretical development (see
\Cref{sec:discussion}). The experimental results are also limited to
relatively small datasets; scaling the diagnostic to large-scale applications
is left for future work.
